\chapter{Particle-Flow Foundations}
\label{ch:bf-dpf-particle-flow-foundations}

This chapter develops the mathematical foundations of particle-flow methods.
Where the previous chapter fixed the classical particle-filter baseline, the
present chapter studies a different idea: instead of relying on repeated
reweighting and resampling to move from the predictive law to the filtering law,
one may attempt to transport particles continuously along a pseudo-time path.

This transport point of view is attractive because it addresses a real weakness
of classical particle filters: in sharply informative or moderately
high-dimensional problems, repeated weight concentration can make the empirical
measure numerically poor long before the state-space model itself ceases to be
well defined.  Particle-flow methods attempt to reduce that instability by
constructing a velocity field that pushes the predictive cloud toward the
posterior cloud.  The mathematical difficulty is that the transport law is not
free: one must define the interpolating densities, the continuity equation, the
velocity field, and the assumptions under which the resulting flow is exact or
only approximate.

This chapter therefore treats particle flow first as a transport problem, not
as a final likelihood engine.  In particular, the present chapter does
\emph{not} yet claim that a raw particle flow defines the final HMC target.  It
first develops the homotopy and flow mathematics, identifies the Gaussian-
closure approximation behind the classical EDH construction, and states the
exact special case in which the transport recovers the Kalman update.  The
subsequent PF-PF chapter then asks how to restore a clearer proposal/target
relationship by importance correction.

\section{Objects, notation, and source roles}
\label{sec:bf-pff-objects}

The particle-flow literature is easy to misread because the same word "flow"
can refer to several distinct objects: a density path, a velocity field, a
particle ODE, a numerical map produced by integrating that ODE, or a proposal
mechanism later corrected by an importance weight.  This chapter uses the
following objects.

\begin{center}
\small
\begin{tabular}{p{0.22\textwidth} p{0.68\textwidth}}
\toprule
Symbol or phrase & Meaning in this chapter \\
\midrule
$\pi_{t|t-1}$ & predictive law of $x_t$ given $y_{1:t-1}$ before assimilating
$y_t$ \\
$\pi_{t|t}$ & filtering law of $x_t$ given $y_{1:t}$ after assimilating
$y_t$ \\
$\pi_{t,\lambda}$ & normalized homotopy density at pseudo-time
$\lambda\in[0,1]$ \\
$Z_t(\lambda)$ & normalizing constant of the homotopy path \\
$v_{t,\lambda}$ & velocity field whose flow is intended to transport
$\pi_{t,\lambda}$ \\
$\Phi_{t,\lambda}$ & flow map generated by the particle ODE from pseudo-time
$0$ to $\lambda$ \\
$A_t(\lambda),b_t(\lambda)$ & affine EDH velocity coefficients under a
Gaussian closure \\
$P_{t,\lambda},m_{t,\lambda}$ & covariance and mean of the Gaussian homotopy
closure \\
$G_t^{(i)}$ & particle-local observation Jacobian in LEDH \\
$\Lambda_t,\Lambda_t^{(i)}$ & global and particle-local likelihood precision
contributions \\
$\eta_t^{(i)}$ & particle-local Gaussian information vector in LEDH \\
\bottomrule
\end{tabular}
\end{center}

The source role is also deliberately limited.  Daum--Huang particle flow
\citep{daumhuang2008} motivates the EDH construction.  Invertible particle-flow
particle filters \citep{li2017particle} explain why later chapters must treat a
flow as a proposal map with a density correction rather than as an automatically
correct posterior sampler.  Particle-flow applications such as
\citet{hu2021particle} motivate numerical and high-dimensional caution.  These
sources locate the method family; they do not remove the need for the local
derivations below, and they do not establish a final HMC target for BayesFilter.

\section{From discrete reweighting to continuous transport}
\label{sec:bf-pff-motivation}

Let the predictive law at time $t$ be
\[
  \pi_{t|t-1}(dx) = p_\theta(x_t \in dx \mid y_{1:t-1}),
\]
and let the filtering law be
\[
  \pi_{t|t}(dx) = p_\theta(x_t \in dx \mid y_{1:t}).
\]
The bootstrap particle filter approximates the transition from
$\pi_{t|t-1}$ to $\pi_{t|t}$ by assigning observation-driven weights and then
resampling the cloud.  Particle-flow methods replace this discrete update by a
continuous pseudo-time transport from a prior-like distribution to a
posterior-like distribution.  The intent is to move particles toward regions of
high posterior density before resampling is needed, thereby reducing weight
collapse.

The conceptual shift is therefore from
\begin{enumerate}[label=(\roman*)]
  \item a \emph{weighted empirical update} followed by discrete resampling,
\end{enumerate}
to
\begin{enumerate}[label=(\roman*),resume]
  \item a \emph{continuous transport path} whose endpoint approximates the
    filtering law.
\end{enumerate}
The advantage of the second view is geometric: the cloud is moved by a velocity
field rather than by repeated duplication and deletion of particles.  The cost
is mathematical and numerical: the transport field must be derived, integrated,
and interpreted correctly.

\section{Homotopy path between predictive and filtering laws}
\label{sec:bf-pff-homotopy}

Particle-flow methods introduce a pseudo-time parameter
$\lambda \in [0,1]$ and construct an interpolating family of densities
$\pi_{t,\lambda}$.  Let
\[
  \bar\pi_{t,\lambda}(x)
  =
  p_\theta(x_t=x \mid y_{1:t-1})
  p_\theta(y_t \mid x)^\lambda
\]
be the unnormalized path.  The normalized path is
\[
  \pi_{t,\lambda}(x)
  =
  \frac{\bar\pi_{t,\lambda}(x)}{Z_t(\lambda)}.
\]
Equivalently, in the classical Daum--Huang construction one defines the
log-homotopy
\begin{equation}
\label{eq:bf-pff-log-homotopy}
  \log \pi_{t,\lambda}(x)
  = \log p_\theta(x_t=x \mid y_{1:t-1})
    + \lambda \log p_\theta(y_t \mid x)
    - \log Z_t(\lambda),
\end{equation}
where
\begin{equation}
\label{eq:bf-pff-homotopy-normalizer}
  Z_t(\lambda)
  = \int p_\theta(x_t \mid y_{1:t-1})
      p_\theta(y_t \mid x_t)^\lambda\,dx_t
\end{equation}
is the normalizing constant.

This definition deserves more emphasis than a short formula often receives.
First, the endpoint interpretation is exact:
\[
  \pi_{t,0}(x) = p_\theta(x_t=x \mid y_{1:t-1}),
  \qquad
  \pi_{t,1}(x) = p_\theta(x_t=x \mid y_{1:t}).
\]
Thus the family joins the predictive law to the filtering law without changing
what the endpoints mean.  Second, the normalization is not merely a technical
constant.  It is the term that keeps the intermediate object a probability
density rather than an unnormalized score.  Third, the homotopy does not yet
choose a transport field; it only chooses the density path the field is supposed
to realize.

Differentiating \eqref{eq:bf-pff-log-homotopy} with respect to $\lambda$ gives
\begin{equation}
\label{eq:bf-pff-log-homotopy-derivative}
  \partial_\lambda \log \pi_{t,\lambda}(x)
  = \log p_\theta(y_t \mid x) - \partial_\lambda \log Z_t(\lambda).
\end{equation}
The term involving $Z_t(\lambda)$ is crucial: it keeps the intermediate density
normalized.  Omitting it may produce a convenient algebraic simplification, but
it changes the law being transported.  This is one of the first places where a
future implementation can quietly go wrong: if the code is written only in terms
of an unnormalized path, the resulting transported object may no longer match
the intended predictive-to-filtering homotopy.

One can make the normalizer term explicit.  Differentiating
\eqref{eq:bf-pff-homotopy-normalizer} gives
\[
  \frac{dZ_t(\lambda)}{d\lambda}
  =
  \int p_\theta(x_t \mid y_{1:t-1})
       p_\theta(y_t \mid x_t)^\lambda
       \log p_\theta(y_t \mid x_t)\,dx_t.
\]
This differentiation under the integral sign is not automatic.  A sufficient
local condition is that, on the pseudo-time interval being used, the likelihood
is positive where the predictive density has support and there is an integrable
envelope dominating
\[
  p_\theta(x_t\mid y_{1:t-1})
  p_\theta(y_t\mid x_t)^\lambda
  \left|\log p_\theta(y_t\mid x_t)\right|.
\]
Equivalently, the predictive law must give finite expectation to the
log-likelihood factor along the homotopy and permit dominated differentiation
of $Z_t(\lambda)$.  If this condition fails, the centered identity below should
be treated as a formal diagnostic rather than as a justified derivative
calculation.
Therefore
\[
  \partial_\lambda \log Z_t(\lambda)
  =
  \int \log p_\theta(y_t \mid x)\,
       \pi_{t,\lambda}(x)\,dx.
\]
The derivative in \eqref{eq:bf-pff-log-homotopy-derivative} is thus the
centered log-likelihood under the intermediate law:
\[
  \partial_\lambda \log \pi_{t,\lambda}(x)
  =
  \log p_\theta(y_t \mid x)
  -
  \mathbb E_{\pi_{t,\lambda}}
  [\log p_\theta(y_t \mid X_t)].
\]
This centered form is a useful audit identity.  It shows that the density path
preserves total probability, since integrating
$\partial_\lambda\pi_{t,\lambda}$ over $x$ gives zero.

For debugging purposes, the main lesson of this section is simple: if a claimed
particle-flow implementation does not recover the correct endpoint target even in
a simple controlled case, one of the first questions to ask is whether the
normalized homotopy was implemented or whether an unnormalized shortcut was
silently substituted.

\section{Continuity equation and transport PDE}
\label{sec:bf-pff-continuity}

Suppose the particle cloud is transported by a velocity field
$v_{t,\lambda}(x)$ through the pseudo-time ODE
\begin{equation}
\label{eq:bf-pff-particle-ode}
  \frac{d x(\lambda)}{d\lambda} = v_{t,\lambda}(x(\lambda)).
\end{equation}
The following derivation is a classical strong-form calculation.  It assumes
that the density is differentiable in $\lambda$, continuously differentiable in
$x$ on the region of interest, that
$\pi_{t,\lambda}v_{t,\lambda}$ is integrable, and that the boundary term in the
integration-by-parts step vanishes.  This is satisfied, for example, for smooth
rapidly decaying densities on $\mathbb R^{n_x}$ with controlled velocity, for
compactly supported smooth test functions in the weak form, or for bounded
domains with an explicitly imposed no-flux or otherwise stated boundary
condition.  Without one of these admissible settings, the PDE should be read as
a formal transport equation rather than a fully justified global statement.
For reviewer-grade use, the weak-form identity with compactly supported test
functions is the safer default; the displayed strong-form PDE is used only when
the stated differentiability, integrability, and boundary-flux hypotheses hold.
Mass conservation over any regular test region then implies that the evolving
density must satisfy the continuity equation
\begin{equation}
\label{eq:bf-pff-continuity-equation}
  \partial_\lambda \pi_{t,\lambda}(x)
  + \nabla \cdot \bigl(\pi_{t,\lambda}(x)
  v_{t,\lambda}(x)\bigr) = 0.
\end{equation}
Equivalently, for any smooth compactly supported test function $\varphi$,
\[
  \frac{d}{d\lambda}
  \int \varphi(x)\pi_{t,\lambda}(x)\,dx
  =
  \int \nabla\varphi(x)^\top
  v_{t,\lambda}(x)\pi_{t,\lambda}(x)\,dx,
\]
and integration by parts gives \eqref{eq:bf-pff-continuity-equation}.
Dividing by $\pi_{t,\lambda}(x)$ where the density is positive yields
\begin{equation}
\label{eq:bf-pff-log-density-continuity}
  \partial_\lambda \log \pi_{t,\lambda}(x)
  = -\nabla \cdot v_{t,\lambda}(x)
    - \nabla \log \pi_{t,\lambda}(x)^\top v_{t,\lambda}(x).
\end{equation}
Combining \eqref{eq:bf-pff-log-homotopy-derivative} and
\eqref{eq:bf-pff-log-density-continuity} gives the basic particle-flow PDE:
\begin{equation}
\label{eq:bf-pff-flow-pde}
  \log p_\theta(y_t \mid x)
  - \partial_\lambda \log Z_t(\lambda)
  = -\nabla \cdot v_{t,\lambda}(x)
    - \nabla \log \pi_{t,\lambda}(x)^\top v_{t,\lambda}(x).
\end{equation}
This equation is exact as a conservation statement.  What is not yet fixed is a
unique velocity field.  The PDE constrains the field but does not determine it
without additional structure.  In particular, if a vector field has zero
weighted divergence with respect to $\pi_{t,\lambda}$, adding it to one
solution can leave the same density evolution.  EDH and LEDH should therefore
be read as particular closure choices, not as the only possible solution of the
transport PDE.

\section{EDH under Gaussian closure}
\label{sec:bf-pff-edh}

The classical exact Daum--Huang (EDH) construction becomes explicit after a
Gaussian closure.  The word "exact" in the name is historical and must be read
with the regime stated here: the derivation is exact for the Gaussian family it
postulates, and it is an exact filtering transport only when that Gaussian
family is the true homotopy family.

Assume that the predictive law is approximated by
\[
  p_\theta(x_t=x \mid y_{1:t-1})
  \approx \mathcal N(x;m_{t|t-1},P_{t|t-1})
\]
and that the observation density is Gaussian with linear observation map
$H_t$ and covariance $R_t$:
\[
  p_\theta(y_t\mid x)
  \propto
  \exp\left\{
    -\frac{1}{2}(y_t-H_tx)^\top R_t^{-1}(y_t-H_tx)
  \right\}.
\]
Here $P_{t|t-1}$ and $R_t$ are assumed symmetric positive definite, and the
matrix products have the conformable dimensions
$H_t\in\mathbb R^{n_y\times n_x}$,
$P_{t|t-1}\in\mathbb R^{n_x\times n_x}$, and
$R_t\in\mathbb R^{n_y\times n_y}$.  These assumptions are what make the inverse
and solve notation below meaningful; an implementation should use linear solves
and record conditioning diagnostics rather than treating the displayed inverses
as a numerical instruction.
At pseudo-time $\lambda$ the interpolating density is then approximated by
\[
  \pi_{t,\lambda}(x) \approx \mathcal N(x; m_{t,\lambda}, P_{t,\lambda}).
\]
The likelihood precision is
\begin{equation}
\label{eq:bf-pff-likelihood-precision}
  \Lambda_t = H_t^\top R_t^{-1} H_t.
\end{equation}
Collecting the quadratic terms in the log-homotopy gives
\[
  P_{t,\lambda}^{-1}
  =
  P_{t|t-1}^{-1}+\lambda\Lambda_t,
\]
so the homotopy covariance satisfies
\begin{equation}
\label{eq:bf-pff-homotopy-covariance}
  P_{t,\lambda} = \bigl(P_{t|t-1}^{-1} + \lambda \Lambda_t\bigr)^{-1},
\end{equation}
and the matrix inverse derivative gives the covariance evolution.  To see the
step explicitly, set
$M_{t,\lambda}=P_{t|t-1}^{-1}+\lambda\Lambda_t$, so that
$M_{t,\lambda}P_{t,\lambda}=I$.  Differentiating this identity yields
$\Lambda_tP_{t,\lambda}+M_{t,\lambda}(dP_{t,\lambda}/d\lambda)=0$ and hence
\begin{equation}
\label{eq:bf-pff-homotopy-cov-derivative}
  \frac{dP_{t,\lambda}}{d\lambda}
  = - P_{t,\lambda} \Lambda_t P_{t,\lambda}.
\end{equation}
The corresponding information vector is
\[
  h_{t,\lambda}
  =
  P_{t|t-1}^{-1}m_{t|t-1}
  + \lambda H_t^\top R_t^{-1}y_t,
  \qquad
  m_{t,\lambda}=P_{t,\lambda}h_{t,\lambda}.
\]
Differentiating $m_{t,\lambda}$ gives
\[
  \frac{dm_{t,\lambda}}{d\lambda}
  =
  -P_{t,\lambda}\Lambda_t m_{t,\lambda}
  + P_{t,\lambda}H_t^\top R_t^{-1}y_t.
\]

The EDH ansatz now chooses an affine velocity field
\begin{equation}
\label{eq:bf-pff-edh-velocity}
  v_{t,\lambda}(x) = A_t(\lambda)x + b_t(\lambda),
\end{equation}
whose mean and covariance evolution match the Gaussian homotopy.  If
$X_\lambda\sim\mathcal N(m_{t,\lambda},P_{t,\lambda})$ evolves under this
affine field, then
\[
  \frac{dm_{t,\lambda}}{d\lambda}
  =
  A_t(\lambda)m_{t,\lambda}+b_t(\lambda),
  \qquad
  \frac{dP_{t,\lambda}}{d\lambda}
  =
  A_t(\lambda)P_{t,\lambda}
  +
  P_{t,\lambda}A_t(\lambda)^\top.
\]
Choosing
$A_t(\lambda)=-\frac{1}{2}P_{t,\lambda}\Lambda_t$ matches
\eqref{eq:bf-pff-homotopy-cov-derivative}.  The offset $b_t(\lambda)$ is then
determined by the mean equation:
\[
  b_t(\lambda)
  =
  \frac{dm_{t,\lambda}}{d\lambda}
  - A_t(\lambda)m_{t,\lambda}.
\]
Using
\[
  m_{t,\lambda}
  =
  m_{t|t-1}
  + \lambda P_{t,\lambda}H_t^\top R_t^{-1}
    (y_t-H_tm_{t|t-1})
\]
rewrites this offset in the standard BayesFilter EDH convention.  The
coefficients are
\begin{align}
\label{eq:bf-pff-edh-A}
  A_t(\lambda)
  &= -\frac{1}{2} P_{t,\lambda}\Lambda_t, \\
\label{eq:bf-pff-edh-b}
  b_t(\lambda)
  &= -A_t(\lambda)m_{t|t-1}
     + \bigl(I + \lambda A_t(\lambda)\bigr)
       P_{t,\lambda}H_t^\top R_t^{-1}
       (y_t - H_t m_{t|t-1}).
\end{align}
Thus the EDH pseudo-time ODE is
\begin{equation}
\label{eq:bf-pff-edh-ode}
  \frac{dx}{d\lambda} = A_t(\lambda)x + b_t(\lambda).
\end{equation}

The Gaussian closure should be read with care.  The continuity equation and the
homotopy construction are exact at their own level.  The approximation enters
when one assumes that the intermediate family remains Gaussian and then chooses
an affine velocity field that is consistent with that Gaussian family.  In other
words, EDH is exact as a transport formula only when the closure assumption is
exactly true; otherwise it is a mathematically structured approximation.

A future implementation should therefore separate four tasks explicitly:
\begin{enumerate}[label=(\roman*)]
  \item evaluate the predictive Gaussian inputs $(m_{t|t-1},P_{t|t-1})$;
  \item build the precision contribution $\Lambda_t$;
  \item build the information contribution $H_t^\top R_t^{-1}y_t$ or its
    locally shifted analogue;
  \item integrate the affine pseudo-time dynamics
    \eqref{eq:bf-pff-edh-ode} using a numerically stable scheme.
\end{enumerate}
If the transported cloud behaves unexpectedly, the first debugging question is
not merely whether the ODE solver was stable, but also whether the Gaussian
closure itself was plausible for the state-space regime under study.

\section{Linear-Gaussian recovery as the exact special case}
\label{sec:bf-pff-linear-gaussian-recovery}

The most important exactness statement for EDH is the linear-Gaussian recovery
result.  Suppose the predictive law is Gaussian and the observation model is
truly linear-Gaussian.  Then the homotopy family remains Gaussian for all
$\lambda$.  The covariance endpoint at $\lambda=1$ is
\begin{equation}
\label{eq:bf-pff-kalman-endpoint}
  P_{t,1} = \bigl(P_{t|t-1}^{-1} + H_t^\top R_t^{-1}H_t\bigr)^{-1},
\end{equation}
which is exactly the Kalman posterior covariance in information form.  The mean
endpoint follows from the endpoint information vector:
\begin{equation}
\label{eq:bf-pff-kalman-mean-endpoint}
  m_{t,1}
  =
  P_{t,1}
  \left(
    P_{t|t-1}^{-1}m_{t|t-1}
    + H_t^\top R_t^{-1}y_t
  \right).
\end{equation}
Using the Kalman gain
\[
  K_t=P_{t|t-1}H_t^\top
  (H_tP_{t|t-1}H_t^\top+R_t)^{-1},
\]
the same mean can be written as
\[
  m_{t,1}
  =
  m_{t|t-1}+K_t(y_t-H_tm_{t|t-1}).
\]
Thus, in the linear-Gaussian case, EDH is not merely an approximation: it
transports the predictive Gaussian law to the exact Kalman update, provided the
affine ODE is integrated exactly or to controlled numerical tolerance.

This statement is crucial for BayesFilter because it supplies the exact special
case against which any EDH implementation should be tested.  It is also the
boundary line for the rest of the chapter: outside this special case, the
particle-flow path is no longer guaranteed exact merely by writing down the
homotopy.

\section{LEDH and local linearization}
\label{sec:bf-pff-ledh}

The local EDH (LEDH) construction replaces the single global observation
linearization by a particle-specific local linearization.  This is not merely
an implementation refinement.  It changes the closure from one Gaussian
homotopy shared by the cloud to a family of particle-local Gaussian
approximations.

For particle $x^{(i)}$, let
\begin{equation}
\label{eq:bf-pff-local-jacobian}
  G_t^{(i)} = \left.\frac{\partial g_\theta(x)}{\partial x}\right|_{x=x^{(i)}}
\end{equation}
and write the local linear approximation of the observation map as
\begin{equation}
\label{eq:bf-pff-local-linearization}
  g_\theta(x)
  \approx g_\theta(x^{(i)}) + G_t^{(i)}(x-x^{(i)}).
\end{equation}
Equivalently, the local observation equation is approximated by
\[
  y_t
  \approx
  G_t^{(i)}x
  +
  \bigl(g_\theta(x^{(i)})-G_t^{(i)}x^{(i)}\bigr)
  +
  \varepsilon_t,
  \qquad
  \varepsilon_t\sim\mathcal N(0,R_t).
\]
Collecting quadratic and linear terms in this local Gaussian likelihood gives
the local precision
\begin{equation}
\label{eq:bf-pff-local-precision}
  \Lambda_t^{(i)} = (G_t^{(i)})^\top R_t^{-1} G_t^{(i)}
\end{equation}
and the local Gaussian information vector
\begin{equation}
\label{eq:bf-pff-local-information-vector}
  \eta_t^{(i)}
  = (G_t^{(i)})^\top R_t^{-1}
    \bigl(y_t - g_\theta(x^{(i)}) + G_t^{(i)}x^{(i)}\bigr),
\end{equation}
which is the particle-wise analogue of the global Gaussian information term
$H_t^\top R_t^{-1}y_t$ in EDH.  The corresponding local homotopy covariance is
\begin{equation}
\label{eq:bf-pff-local-homotopy-covariance}
  P_{t,\lambda}^{(i)}
  = \bigl(P_{t|t-1}^{-1} + \lambda \Lambda_t^{(i)}\bigr)^{-1}.
\end{equation}
The associated local mean is
\[
  m_{t,\lambda}^{(i)}
  =
  P_{t,\lambda}^{(i)}
  \left(P_{t|t-1}^{-1}m_{t|t-1}
  +\lambda\eta_t^{(i)}\right).
\]
The local linearization of the observation map near particle $x^{(i)}$ therefore
produces particle-specific flow coefficients
\begin{align}
\label{eq:bf-pff-ledh-A}
  A_t^{(i)}(\lambda)
  &= -\frac{1}{2} P_{t,\lambda}^{(i)} \Lambda_t^{(i)}, \\
\label{eq:bf-pff-ledh-b}
  b_t^{(i)}(\lambda)
  &= -A_t^{(i)}(\lambda)m_{t|t-1}
     + \bigl(I + \lambda A_t^{(i)}(\lambda)\bigr)
       P_{t,\lambda}^{(i)} \eta_t^{(i)},
\end{align}
with pseudo-time evolution
\begin{equation}
\label{eq:bf-pff-ledh-ode}
  \frac{dx^{(i)}}{d\lambda}
  = A_t^{(i)}(\lambda)x^{(i)} + b_t^{(i)}(\lambda).
\end{equation}

For the Li--Coates PF-PF(LEDH) construction used in the next chapter,
there is an additional implementation-facing distinction that must not be
collapsed.  The local linearization is not a single deterministic affine map
shared by all particles, and it is not computed from a cloud-level average.
For particle $i$, Algorithm~1 of \citet{li2017particle} first builds a
particle-specific predicted covariance $P^i$ and a zero-noise transition anchor
\[
  \bar\eta_0^i = g_k(x_{k-1}^i,0).
\]
It also samples an actual pre-flow proposal particle
\[
  \eta_0^i = g_k(x_{k-1}^i,v_k).
\]
During pseudo-time integration, the Jacobian used in the LEDH coefficients is
evaluated at the moving auxiliary state $\bar\eta^i_\lambda$, not at a fixed
global mean and not at another particle.  In the notation of
\citet{li2017particle}, the local observation linearization is
\begin{equation}
\label{eq:bf-pff-li-coates-local-jacobian}
  H^i(\lambda)
  =
  \left.
  \frac{\partial h(\eta,0)}{\partial \eta}
  \right|_{\eta=\bar\eta^i_\lambda},
  \qquad
  e^i(\lambda)
  =
  h(\bar\eta^i_\lambda,0)-H^i(\lambda)\bar\eta^i_\lambda .
\end{equation}
With the particle-specific prediction covariance $P^i$ held as the covariance
input for that particle's flow step, the source-form LEDH coefficients are
\begin{align}
\label{eq:bf-pff-li-coates-ledh-A}
  A^i(\lambda)
  &=
  -\frac{1}{2}
  P^i H^i(\lambda)^\top
  \bigl(\lambda H^i(\lambda)P^iH^i(\lambda)^\top+R\bigr)^{-1}
  H^i(\lambda), \\
\label{eq:bf-pff-li-coates-ledh-b}
  b^i(\lambda)
  &=
  \bigl(I+2\lambda A^i(\lambda)\bigr)
  \left[
    \bigl(I+\lambda A^i(\lambda)\bigr)
    P^i H^i(\lambda)^\top R^{-1}\{z_k-e^i(\lambda)\}
    + A^i(\lambda)\bar\eta_0^i
  \right].
\end{align}
The auxiliary anchor and the actual proposal particle are then migrated by the
same particle-local affine update at each pseudo-time step:
\begin{align}
\label{eq:bf-pff-li-coates-aux-update}
  \bar\eta^i_{\lambda_j}
  &=
  \bar\eta^i_{\lambda_{j-1}}
  +\epsilon_j
  \left[
    A_j^i(\lambda_j)\bar\eta^i_{\lambda_{j-1}}
    +b_j^i(\lambda_j)
  \right],\\
\label{eq:bf-pff-li-coates-actual-update}
  \eta^i_{\lambda_j}
  &=
  \eta^i_{\lambda_{j-1}}
  +\epsilon_j
  \left[
    A_j^i(\lambda_j)\eta^i_{\lambda_{j-1}}
    +b_j^i(\lambda_j)
  \right].
\end{align}
This auxiliary/actual split is the key source-faithfulness point for LEDH in
PF-PF.  The auxiliary path determines the local linearization, while the actual
path is the transported proposal particle whose density must later be corrected.
Using a single global affine map, or linearizing every particle at a fixed
cloud-level point, is therefore not Li--Coates Algorithm~1 LEDH.

The same source-faithfulness point applies to the per-particle covariance
objects.  The covariance $P^i$ is part of the Li--Coates Algorithm~1 state used
to construct the particle-local LEDH coefficients; BayesFilter does not treat
this covariance state as a new filter invention.  When a later differentiable
OT or Sinkhorn resampling layer is attached to this Algorithm~1 route, the
additional BayesFilter obligation is bookkeeping: specify how the covariance
auxiliary state is carried through the chosen transport map, and test that this
carry rule preserves the intended route.  The transport layer supplies the
resampling relaxation; it does not replace the source algorithm's covariance
lifecycle or by itself prove HMC suitability.

This construction explains both the attraction and the danger of LEDH.  The
attraction is that each particle uses a local approximation intended to track
nonlinear observation geometry more closely than a single global
linearization.  The danger is that the transported law now depends on a family
of particle-specific linearizations, each of which may behave poorly if the
Jacobian is unstable, ill-conditioned, or evaluated in a region where the local
approximation is misleading.  Hence LEDH is potentially more adaptive than EDH,
but it is also more fragile numerically and interpretively.  The chapter should
not phrase LEDH as guaranteed improvement without model-specific evidence.

From an implementation and debugging standpoint, LEDH adds at least three new
failure routes relative to EDH:
\begin{enumerate}[label=(\roman*)]
  \item inaccurate or unstable local Jacobians;
  \item poor local-information-vector construction;
  \item path-dependent variation in stiffness across particles.
\end{enumerate}
These are not secondary engineering details.  They are direct consequences of
how the local approximation enters the mathematics.

\section{Stiffness and discretization}
\label{sec:bf-pff-stiffness}

Even when the flow equations are well defined, their numerical integration may
be difficult.  The main difficulty is stiffness.  When the observation noise is
small or the effective likelihood precision is large, the eigenvalues of the
drift matrix in the EDH or LEDH ODE can span several orders of magnitude.  In
that regime, explicit solvers require very small pseudo-time steps to avoid
large truncation error or instability.

Formally, the stiffness is tied to the conditioning of the homotopy covariance
or precision family.  In the local case, for example, large eigenvalue spreads
in $\Lambda_t^{(i)}$ may force extremely fine pseudo-time resolution near the
posterior end of the path.  This matters for implementation because the flow
construction is only useful if it can be integrated stably and repeatedly inside
a filtering recursion.

The practical burden can be stated more concretely.  If the pseudo-time ODE is
integrated with an explicit method such as RK4, then the local truncation error
is governed by powers of the step size, but the stability envelope is controlled
by the operator norm and spectral spread of the drift matrix.  In sharp-
likelihood regimes, decreasing the observation noise or increasing the local
precision can force the stable step size down dramatically.  For LEDH this may
happen non-uniformly across the cloud, because different particles see different
local Jacobians and therefore different local precision matrices.

This is exactly why stiffness belongs in the monograph as more than a short
warning.  It creates a direct literature-to-debugging bridge:
\begin{itemize}
  \item if a flow implementation behaves erratically only when observation noise
    becomes small, stiffness is a prime mathematical suspect;
  \item if certain particles show unstable trajectories while others remain well
    behaved, the local LEDH precision field is a prime suspect;
  \item if a nominally correct flow map fails under a coarse pseudo-time grid,
    the first question should be whether the flow is underresolved rather than
    whether the underlying theory is wrong.
\end{itemize}

For BayesFilter, the main consequence is therefore both conceptual and
operational.  Stiffness is not merely a performance issue.  It is one of the
main ways in which a formally attractive flow construction can become
numerically unreliable enough to alter the practical behavior of the filter.
That is why later implementation-facing chapters must record discretization
policy, stiffness diagnostics, and stress tests explicitly, and why the
literature pointers for stiffness should be treated as debugging tools rather
than as optional background reading.

\section{Exactness and approximation ledger}
\label{sec:bf-pff-claim-ledger}

The particle-flow chapter has several layers that must not be collapsed.

\begin{center}
\small
\begin{tabular}{p{0.23\textwidth} p{0.34\textwidth} p{0.32\textwidth}}
\toprule
Layer & Status & Main implementation diagnostic \\
\midrule
Normalized homotopy endpoints & exact model identity when the predictive
density and likelihood are the intended model objects & recover predictive law
at $\lambda=0$ and filtering law at $\lambda=1$ in analytic test cases \\
Continuity equation & exact conservation statement under regularity and
boundary assumptions & verify transported test-function moments against ODE
integration in controlled problems \\
Velocity field choice & not unique; EDH/LEDH are closure choices & record the
chosen field and compare alternatives on benchmark models \\
EDH Gaussian closure & exact for the Gaussian family; exact filtering only in
linear-Gaussian recovery & test covariance and mean endpoints against the
Kalman update \\
LEDH local linearization & approximate closure using particle-local Jacobians
and information vectors & monitor Jacobian conditioning, local precision
spectra, and particle-wise endpoint dispersion \\
Pseudo-time integration & numerical approximation to the chosen ODE & run step
size, tolerance, stiffness, and log-det sensitivity checks \\
\bottomrule
\end{tabular}
\end{center}

This ledger is deliberately conservative.  It permits a strong statement about
linear-Gaussian recovery, but it blocks any inference that a raw EDH or LEDH
cloud is already a corrected nonlinear filter, a pseudo-marginal likelihood
estimator, or an HMC-ready target.

\section{What this chapter does and does not claim}
\label{sec:bf-pff-chapter-boundary}

This chapter develops the particle-flow transport mathematics itself.  It does
\emph{not} yet claim:
\begin{itemize}
  \item that the transported cloud alone defines the final corrected filtering
    law in the nonlinear case;
  \item that raw EDH or LEDH transport is already the final HMC target;
  \item that the Jacobian/proposal correction issue has been resolved;
  \item that differentiable resampling or OT relaxations have been introduced.
\end{itemize}
Those belong to the later PF-PF and resampling chapters.  The role of the
present chapter is narrower and mathematically prior: it explains how the
transport path is constructed, where the Gaussian-closure or local-linearization
approximation enters, and why the linear-Gaussian case supplies the exact
special benchmark.
